%% 
%% Copyright 2007-2024 Elsevier Ltd
%% 
%% This file is part of the 'Elsarticle Bundle'.
%% ---------------------------------------------
%% 
%% It may be distributed under the conditions of the LaTeX Project Public
%% License, either version 1.3 of this license or (at your option) any
%% later version.  The latest version of this license is in
%%    http://www.latex-project.org/lppl.txt
%% and version 1.3 or later is part of all distributions of LaTeX
%% version 1999/12/01 or later.
%% 
%% The list of all files belonging to the 'Elsarticle Bundle' is
%% given in the file `manifest.txt'.
%% 
%% Template article for Elsevier's document class `elsarticle'
%% with harvard style bibliographic references

\documentclass[preprint,12pt,authoryear]{elsarticle}

%% Use the option review to obtain double line spacing
%% \documentclass[authoryear,preprint,review,12pt]{elsarticle}

%% Use the options 1p,twocolumn; 3p; 3p,twocolumn; 5p; or 5p,twocolumn
%% for a journal layout:
%% \documentclass[final,1p,times,authoryear]{elsarticle}
%% \documentclass[final,1p,times,twocolumn,authoryear]{elsarticle}
%% \documentclass[final,3p,times,authoryear]{elsarticle}
%% \documentclass[final,3p,times,twocolumn,authoryear]{elsarticle}
%% \documentclass[final,5p,times,authoryear]{elsarticle}
%% \documentclass[final,5p,times,twocolumn,authoryear]{elsarticle}

%% For including figures, graphicx.sty has been loaded in
%% elsarticle.cls. If you prefer to use the old commands
%% please give \usepackage{epsfig}

%% The amssymb package provides various useful mathematical symbols
\usepackage{amssymb}
%% The amsmath package provides various useful equation environments.
\usepackage{amsmath}
%% The amsthm package provides extended theorem environments
%% \usepackage{amsthm}

%% The lineno packages adds line numbers. Start line numbering with
%% \begin{linenumbers}, end it with \end{linenumbers}. Or switch it on
%% for the whole article with \linenumbers.
%% \usepackage{lineno}

\journal{Smart and Sustainable Built Environment?}

\begin{document}

\begin{frontmatter}

%% Title, authors and addresses

%% use the tnoteref command within \title for footnotes;
%% use the tnotetext command for theassociated footnote;
%% use the fnref command within \author or \affiliation for footnotes;
%% use the fntext command for theassociated footnote;
%% use the corref command within \author for corresponding author footnotes;
%% use the cortext command for theassociated footnote;
%% use the ead command for the email address,
%% and the form \ead[url] for the home page:
%% \title{Title\tnoteref{label1}}
%% \tnotetext[label1]{}
%% \author{Name\corref{cor1}\fnref{label2}}
%% \ead{email address}
%% \ead[url]{home page}
%% \fntext[label2]{}
%% \cortext[cor1]{}
%% \affiliation{organization={},
%%            addressline={}, 
%%            city={},
%%            postcode={}, 
%%            state={},
%%            country={}}
%% \fntext[label3]{}

\title{Using co-design workshops to explore community needs and future directions for shared electric mobility hubs} %% Article title

%% use optional labels to link authors explicitly to addresses:
%% \author[label1,label2]{}
%% \affiliation[label1]{organization={},
%%             addressline={},
%%             city={},
%%             postcode={},
%%             state={},
%%             country={}}
%%
%% \affiliation[label2]{organization={},
%%             addressline={},
%%             city={},
%%             postcode={},
%%             state={},
%%             country={}}



\author[TCD]{Brian Caulfield}
\author[TCD]{Margaret O’Mahony}
\author[TCD]{Sarbast Molsem}
\author[ATU_Galway]{Amaya Vega}
\author[ATU_Sligo]{Brian McCann}
\author[ATU_Galway]{James Reynolds} %% Author name

%% Author affiliation
\affiliation[ATU_Galway]{organization={Atlantic Technological University, Galway City campus},%Department and Organization
            addressline={Dublin Road}, 
            city={Galway City},
            postcode={H91 T8NW}, 
            country={Ireland}}

\affiliation[TCD]{organization={Centre for Transport Research, Dept of Civil, Structural and Environmental Engineering, Trinity College Dublin},%Department and Organization
            addressline={University of Dublin}, 
            country={Ireland}}

\affiliation[ATU_Sligo]{organization={Atlantic Technological University, Sligo campus},
            addressline={Ash Lane},
            city={Sligo},
            postcode={F91 YW50}, 
            country={Ireland}}



%% Abstract
\begin{abstract}
Railway stations, bus stops, post offices and similar facilities can often provide an important focus for communities, including as a central hub for the movement of people and goods, and as a place of social interaction. Recent research directions have emphasised the ‘Place’ function of road (and other transport) facilities as a vital and additional dimension alongside the more traditional focus of transport systems for through movement and more local access to adjoining land uses.  Shared e-mobility hubs are a relatively new development in transport, and seek to provide a central location at which travellers and nearby residents might have access to electric cars, e-cargo bikes, e-bikes or other modes that might support more sustainable travel choices, especially in contexts where the provision of high frequency transit services may be challenging (e.g. regional towns, rural communities, low-density suburbs). However, it is unclear form such shared e-mobility hubs should take, and the extent to which the co-location of other (e.g. library, café, bus stop, community house) might enhance usage, safety and security, and local communities.  Four shared e-mobility hubs have recently been implemented in Letterkenny, Galway, Waterford and the suburbs of Dublin. These include e-bikes, e-cargo bikes, and shared electric cars as part of the ‘base’ offering, but it is unclear what additional mobility, community-focused, commercial or other offerings might be desired by communities, or might help to support the use of shared e-mobility hub facilities to improve transport sustatinability. This paper reports the findings from co-design creative workshops held in Letterkenny, Galway, Waterford and Dublin to explore community attitudes and concepts as to how the ‘base’ shared e-mobility hub might be enhanced so as to better meet each community’s transport and other needs. The cross-case comparisons reported in this paper might provide policy-makers, practitioners and researchers further insights into the types of offerings that might be provided at shared e-mobility hubs and other transport facilities in a variety of contexts (small town, secondary town centre, suburban locations etc.) so as to enhance the sustainable mobility and community outcomes of transport facility development.
\end{abstract}

%%Graphical abstract
%\begin{graphicalabstract}
%\includegraphics{grabs}
%\end{graphicalabstract}

%%Research highlights
%\begin{highlights}
%\item Research highlight 1
%\item Research highlight 2
%\end{highlights}

%% Keywords
\begin{keyword}
%% keywords here, in the form: keyword \sep keyword

%% PACS codes here, in the form: \PACS code \sep code

%% MSC codes here, in the form: \MSC code \sep code
%% or \MSC[2008] code \sep code (2000 is the default)

\end{keyword}

\end{frontmatter}

%% Add \usepackage{lineno} before \begin{document} and uncomment 
%% following line to enable line numbers
%% \linenumbers

%% main text
%%

%% Use \section commands to start a section
\section{Introduction}
\label{sec1}
%% Labels are used to cross-reference an item using \ref command.



%% The Appendices part is started with the command \appendix;
%% appendix sections are then done as normal sections
\appendix
\section{Example Appendix Section}
\label{app1}

Appendix text.

%% For citations use: 
%%       \citet{<label>} ==> Lamport (1994)
%%       \citep{<label>} ==> (Lamport, 1994)
%%
Example citation, See \citet{lamport94}.

%% If you have bib database file and want bibtex to generate the
%% bibitems, please use
%%
%%  \bibliographystyle{elsarticle-harv} 
%%  \bibliography{<your bibdatabase>}

%% else use the following coding to input the bibitems directly in the
%% TeX file.

%% Refer following link for more details about bibliography and citations.
%% https://en.wikibooks.org/wiki/LaTeX/Bibliography_Management

\bibliography{references}
\bibliographystyle{plainnat}
%\begin{thebibliography}{00}

%% For authoryear reference style
%% \bibitem[Author(year)]{label}
%% Text of bibliographic item

%\bibitem[Lamport(1994)]{lamport94}
 % Leslie Lamport,
 % \textit{\LaTeX: a document preparation system},
%  Addison Wesley, Massachusetts,
%  2nd edition,
%  1994.


%\bibitem[Currie(2018)]{Currie:2018aa}
%	Graham Currie, 
%	Lies, Damned Lies, AVs, Shared Mobility, and Urban Transit Futures, 
%	2018.



%\end{thebibliography}
\end{document}

\endinput
%%
%% End of file `elsarticle-template-harv.tex'.


